% Part: methods
% Chapter: proofs
% Section: introduction

\documentclass[../../../include/open-logic-section]{subfiles}

\begin{document}

\olfileid{mth}{prf}{int}

\olsection{प्रस्तावना}

प्रास्ताविक तर्कशास्त्राच्या अभ्यासामुळे तुम्हाला एखादी !!a{derivation} पद्धती परिचित असेल---बहुधा नैसर्गिक निगमनाची किंवा फिच शैलीतील !!{derivation} पद्धती, किंवा कदाचित सिद्धता-वृक्षांची पद्धती. या पद्धतींत !!a{formula} सिद्ध करणे किंवा दिलेला युक्तिवाद युक्त आहे हे दाखवणे, अशा सिद्धता तुम्ही केल्याचे आठवत असेल. त्यासाठी अपेक्षित अंतिम निष्कर्ष मिळेपर्यंत तुम्ही त्या पद्धतीचे नियम वापरले होते. तर्कशास्त्रा\emph{विषयी} विचार करतानाही आपण गोष्टी सिद्ध करतो; पण बहुतेक वेळा !!a{derivation} पद्धती वापरत नाही. आपण पाहणार असलेल्या बहुतेक सिद्धता संपूर्णपणे प्रथम-क्रम तर्कशास्त्राच्या भाषेत लिहिलेल्या नसून नैसर्गिक भाषेत असतात; या आवृत्तीत त्या मराठीत आहेत आणि त्यांत काही सांकेतिक भाषाही येते. अशा सिद्धता रचताना सुरुवातीला तुम्हाला प्रश्न पडू शकतात: !!a{derivation} पद्धतीशिवाय एखादी गोष्ट कशी सिद्ध करायची? सुरुवात कुठून करायची? माझी सिद्धता बरोबर आहे हे कसे ठरवायचे?

सिद्धता करण्याचा प्रयत्न सुरू करण्यापूर्वी सिद्धता म्हणजे काय आणि ती कशी रचायची, हे समजणे महत्त्वाचे आहे. नावावरूनच सूचित होते तसे, एखादी गोष्ट सत्य आहे हे दाखवणे हा \emph{सिद्धते}चा उद्देश असतो. संवादाच्या रूपात याचा विचार करता येईल: कोणी तुम्हाला विचारते की एखादी गोष्ट खरी आहे का---उदाहरणार्थ, दोन वगळता प्रत्येक अविभाज्य संख्या विषम आहे का. केवळ “होय” असे उत्तर देणे पुरेसे नाही; त्या व्यक्तीला \emph{का}, हेही जाणून घ्यायचे असेल. अशा वेळी तुम्ही सिद्धता द्याल.

दैनंदिन संभाषणात उत्तराची ढोबळ दिशा सांगणे किंवा अपूर्ण उत्तर देणे पुरेसे ठरू शकते. पण तर्कशास्त्रात आणि गणितात आपल्याला काटेकोर सिद्धता हवी असते: एखादी गोष्ट खरी असल्याविषयी \emph{कोणत्याही} शंकेला जागा राहू नये, असे दाखवायचे असते. याचा अर्थ सिद्धतेतील प्रत्येक पायरीचे समर्थन केले पाहिजे आणि ते समर्थन योग्य असले पाहिजे. उदाहरणार्थ, तुम्ही वापरत असलेले गृहीतक हे सिद्ध करावयाच्या प्रमेयाच्या विधानात खरोखर गृहीत धरलेले असले पाहिजे; व्याख्यांचा वापर अचूक असला पाहिजे; समर्थनासाठी वापरलेली अनुमाने योग्य असली पाहिजेत; इत्यादी.

सहसा आपण एखादे विधान सिद्ध करत असतो. सिद्ध करावयाच्या विधानांना आपण विविध नावे देतो: विधान, प्रमेय, पूर्वप्रमेय किंवा अनुप्रमेय. विधान म्हणजे सिद्ध करावे असे मूलभूत म्हणणे: नोंदवण्याइतके महत्त्वाचे, पण कदाचित फार सखोल नसलेले किंवा वारंवार न वापरले जाणारे. प्रमेय म्हणजे लक्षणीय, महत्त्वाचे विधान. त्याची सिद्धता अनेकदा काही पायऱ्यांत विभागलेली असते. कधी त्याला ते प्रथम सिद्ध करणाऱ्या व्यक्तीचे नाव दिले जाते---उदाहरणार्थ, कँटरचे प्रमेय किंवा लोव्हेनहाइम--स्कोलेम प्रमेय (L\"owenheim--Skolem)---तर कधी त्याच्या विषयावरून नाव दिले जाते---उदाहरणार्थ, संपूर्णता प्रमेय. पूर्वप्रमेय म्हणजे अधिक महत्त्वाचा निष्कर्ष सिद्ध करण्यासाठी वापरले जाणारे विधान किंवा प्रमेय. नावामुळे गोंधळ होऊ शकतो: कधी पूर्वप्रमेये स्वतःही महत्त्वाचे निष्कर्ष असतात आणि ती मांडणाऱ्या व्यक्तीच्या नावाने ओळखली जातात---उदाहरणार्थ, झोर्नचे पूर्वप्रमेय. अनुप्रमेय म्हणजे दुसऱ्या निष्कर्षावरून सहज मिळणारा निष्कर्ष.

सिद्ध करावयाच्या विधानात अनेकदा काही गृहीतके असतात. आपण कोणत्या प्रकारच्या गोष्टींविषयी सिद्धता करतो, हे त्यांमुळे स्पष्ट होते. विधानाची सुरुवात “$!A$ हे $!B \lif !C$ या रूपाचे !!a{formula} असू द्या” किंवा “$\Gamma \Proves !A$ असे समजा”, किंवा अशाच शब्दांनी होऊ शकते. ही त्या विधानाची, प्रमेयाची किंवा पूर्वप्रमेयाची \emph{गृहीतके} आहेत; सिद्धतेत ती सत्य आहेत असे धरता येते. काय सिद्ध करायचे आहे याची व्याप्ती ती निश्चित करतात आणि आपण ज्या वस्तूंविषयी बोलतो त्यांच्यासाठी काही नावेही देतात. उदाहरणार्थ, तुमच्या विधानाची सुरुवात “$!A$ हे $!B \lif !C$ या रूपाचे !!a{formula} असू द्या” अशी असेल, तर तुम्ही केवळ एका विशिष्ट प्रकारच्या सर्व सूत्रांविषयी---म्हणजे सशर्त सूत्रांविषयी---काहीतरी सिद्ध करत आहात. तुमची सिद्धता ज्या $!B \lif !C$ विषयी बोलते ते त्या प्रकारचे स्वेच्छ सशर्त सूत्र आहे, असे समजले जाते.

\end{document}
